\documentclass{article}
\usepackage[utf8]{inputenc}
\usepackage{amsmath}
\usepackage{graphicx}
\usepackage{geometry}
\geometry{a4paper, margin=1in}

\title{Deep Generative Models - Assignment 2}
\author{Taha Majlesi}
\date{November 2025}

\begin{document}

\maketitle

\tableofcontents

\section{Question 1: Normalizing Flow}

\subsection{Part 1: Theoretical Section}

\subsubsection{Sub-part 1: Change of Variables Formula}

\paragraph{Problem 1.1: $Z \sim U(0, 1), X = 2Z + 1$}
Given the transformation $X = 2Z + 1$, the inverse is $Z = \frac{X - 1}{2}$. The Jacobian of the inverse transformation is $\frac{dZ}{dX} = \frac{1}{2}$. The probability density function of $X$ is given by:
$$p_X(x) = p_Z\left(\frac{x-1}{2}\right) \left| \frac{dZ}{dX} \right| = 1 \cdot \frac{1}{2} = \frac{1}{2}$$
The support for $X$ is $[1, 3]$. So, $p_X(x) = \frac{1}{2}$ for $x \in [1, 3]$ and 0 otherwise.

\paragraph{Problem 1.2: $Z \sim U(0, 2), X = \exp(Z)$}
The inverse transformation is $Z = \ln(X)$, with Jacobian $\frac{dZ}{dX} = \frac{1}{X}$. The PDF of $X$ is:
$$p_X(x) = p_Z(\ln(x)) \left| \frac{1}{x} \right| = \frac{1}{2} \cdot \frac{1}{x} = \frac{1}{2x}$$
The support for $X$ is $[e^0, e^2] = [1, e^2]$. So, $p_X(x) = \frac{1}{2x}$ for $x \in [1, e^2]$ and 0 otherwise.

\subsubsection{Sub-part 2: Jacobian Determinant for Coupling Layers}

\paragraph{Problem 2.1: Additive Coupling Layer}
The transformation is $z_1 = x_1$ and $z_2 = x_2 + m(x_1)$. The Jacobian matrix is:
$$ J = \begin{pmatrix} 1 & 0 \\ m'(x_1) & 1 \end{pmatrix} $$
The determinant is $\det(J) = 1$.

\paragraph{Problem 2.2: Affine Coupling Layer}
The transformation is $z_1 = x_1$ and $z_2 = s(x_1)x_2 + t(x_1)$. The Jacobian matrix is:
$$ J = \begin{pmatrix} 1 & 0 \\ s'(x_1)x_2 + t'(x_1) & s(x_1) \end{pmatrix} $$
The determinant is $\det(J) = s(x_1)$.

\subsection{Part 2: Implementation Section}

\subsubsection{MAF Training and Image Generation}
The MAF model was trained on the capsule dataset. The training loss is shown in Figure \ref{fig:maf_loss}.
\begin{figure}[h!]
    \centering
    \includegraphics[width=0.8\textwidth]{images/maf_loss.png}
    \caption{MAF Training Loss}
    \label{fig:maf_loss}
\end{figure}

Generated images from the trained MAF model are shown in Figure \ref{fig:maf_generated}.
\begin{figure}[h!]
    \centering
    \includegraphics[width=0.8\textwidth]{images/maf_generated.png}
    \caption{Images generated by MAF}
    \label{fig:maf_generated}
\end{figure}

\subsubsection{Anomaly Detection}
The anomaly detection performance was evaluated using AUROC. The ROC curve and score distributions are shown in Figures \ref{fig:roc_curve} and \ref{fig:score_dist}.
\begin{figure}[h!]
    \centering
    \includegraphics[width=0.8\textwidth]{images/roc_curve.png}
    \caption{ROC Curve for Anomaly Detection}
    \label{fig:roc_curve}
\end{figure}
\begin{figure}[h!]
    \centering
    \includegraphics[width=0.8\textwidth]{images/score_dist.png}
    \caption{Distribution of Anomaly Scores}
    \label{fig:score_dist}
\end{figure}

\section{Question 2: CycleGAN}

\subsection{Part 1: Theoretical Questions}
The theoretical questions regarding CycleGAN's loss function, architecture, and limitations were answered in the notebook.

\subsection{Part 2: Implementation Section}
The CycleGAN was trained on the horse2zebra dataset. The training losses are shown in Figure \ref{fig:cyclegan_loss}.
\begin{figure}[h!]
    \centering
    \includegraphics[width=0.8\textwidth]{images/cyclegan_loss.png}
    \caption{CycleGAN Training Losses}
    \label{fig:cyclegan_loss}
\end{figure}

Sample translations are shown in Figures \ref{fig:cyclegan_A2B} and \ref{fig:cyclegan_B2A}.
\begin{figure}[h!]
    \centering
    \includegraphics[width=\textwidth]{images/cyclegan_A2B.png}
    \caption{CycleGAN Translation from Domain A to B}
    \label{fig:cyclegan_A2B}
\end{figure}
\begin{figure}[h!]
    \centering
    \includegraphics[width=\textwidth]{images/cyclegan_B2A.png}
    \caption{CycleGAN Translation from Domain B to A}
    \label{fig:cyclegan_B2A}
\end{figure}

\end{document}
